% =====================================================================
%  Section 4 — Experiments
%
%  Drop-in LaTeX for the experimental section. Designed to be
%  % =====================================================================
%  Section 4 — Experiments
%
%  Drop-in LaTeX for the experimental section. Designed to be
%  % =====================================================================
%  Section 4 — Experiments
%
%  Drop-in LaTeX for the experimental section. Designed to be
%  % =====================================================================
%  Section 4 — Experiments
%
%  Drop-in LaTeX for the experimental section. Designed to be
%  \input{paper_draft/section_4_body.tex} after Section 3.
%
%  Numbers in Tables 3 / 4 / 5 and Figure \ref{fig:k-saturation}
%  come from local 5-seed paired experiments in phase1_distillation/:
%
%    drive_q1_per_iter_vs_sample_once_results.json   (Tab. 3 = Q1)
%    tab4_compiled.json                              (Tab. 4 = Q2)
%    drive_pate_K_saturation_results.json            (Tab. 5 + Fig = Q3)
%    drive_pate_canal_combined_results.json          (Tab. 4 PATE+CANAL row)
%
%  All numbers in the tables can be re-extracted by running
%      python3 phase1_distillation/compile_tab4.py
% =====================================================================

\section{Experiments}
\label{sec:experiments}

We evaluate the framework of Section~\ref{sec:privacy} on retinal vessel
segmentation. Three questions structure the section:

\begin{description}
\item[(Q1)] At a \emph{matched user-level} $\varepsilon$, does the
sample-once release model preserve utility relative to the per-iteration
release model that prior work implicitly uses
(\S\ref{ssec:q1-sample-once})?
\item[(Q2)] Does the channel-importance allocation (CANAL) of
Section~\ref{ssec:wf} improve Dice over uniform noise on real DRIVE
features, and does the honest multi-teacher PATE aggregation of
Section~\ref{ssec:pate} compound with it
(\S\ref{ssec:q2-mechanism})?
\item[(Q3)] How does utility scale with the number of teachers $K$, and
where does it saturate (\S\ref{ssec:q3-k-sweep})?
\end{description}

We additionally include a paired-comparison ablation that tests
Corollary~\ref{cor:r14}'s prediction that CANAL's empirical gain on top
of PATE is bounded by $R^{1/4}$ and therefore negligible on RGB fundus
(\S\ref{ssec:q2-mechanism}, last row of Table~\ref{tab:mech}).

% ---------------------------------------------------------------------
\subsection{Experimental setup}
\label{ssec:setup}

\paragraph{Dataset.} We use DRIVE~\cite{staal2004drive}, the
canonical retinal vessel segmentation benchmark: $40$ RGB fundus images
of size $584\!\times\!565$ with expert vessel masks, split into the
canonical $20$ training and $20$ test images. Each training image is one
``user'' for the purposes of user-level DP, consistent with the
sample-once-per-image model of \S\ref{ssec:threat}.

\paragraph{Teacher and student.} Teacher and student are 3-stage
U-Net~\cite{ronneberger2015unet} variants at $96\!\times\!96$ input
resolution. The teacher has base width 32 (bottleneck dimension
$C\!=\!128$); the student has base width 16 (bottleneck dimension
$64$), with a $1\!\times\!1$ adapter projecting student features into
the teacher channel space for the feature-matching loss. Both models
are trained for $60$ epochs with Adam (lr $=\!10^{-3}$). The student
training loop runs for $40$ epochs against the sample-once cached
targets, using a combined CE+Dice task loss against the ground-truth
mask plus an MSE feature-matching loss with weight $\lambda\!=\!0.4$.

\paragraph{Privacy configuration.} We work in zCDP~\cite{bun2016concentrated}
with $\delta\!=\!10^{-5}$ throughout and report user-level
$\varepsilon\!\in\!\{2, 4, 8, 16\}$ via Eq.~\eqref{eq:eps-to-rho}.
Sensitivities follow the per-channel $L_2$ normalisation of
\S\ref{ssec:accounting}: $\Delta_c\!=\!2/K$, with $K\!=\!1$ unless
otherwise noted. The per-iteration model is parameterised by
$T\!=\!40000$, $B\!=\!4$, $N\!=\!20$, giving $Q\!=\!TB/N\!=\!8000$
per-patient releases under that threat model (Table~\ref{tab:privacy}).

\paragraph{Reporting.} Each cell of every table is the mean and standard
deviation of best validation vessel-Dice over $5$ student-init seeds
$\{100, 200, 300, 400, 500\}$. Cross-mechanism lifts are reported as
\emph{paired} per-seed differences (each seed pair shares student
initialisation, dataloader ordering, and the noise sample produced by
the sample-once precompute) with their standard error of the mean. We
report $t$-statistics on the paired lift; a $|t|\!\geq\!2$ corresponds
to $p\!<\!0.05$ in a paired one-sample test.

% ---------------------------------------------------------------------
\subsection{Q1 — sample-once preserves utility at matched user-level $\varepsilon$}
\label{ssec:q1-sample-once}

Table~\ref{tab:q1-sample-once} reports best validation vessel-Dice under
the two release models at every user-level $\varepsilon\!\in\!\{2,4,8,16\}$
that we test. The per-iter row uses the same student training loop as
sample-once, but at every training iteration adds a \emph{fresh} draw of
Gaussian noise on top of the clean teacher bottleneck of the current
batch, scaled so that the cumulative zCDP cost over $Q\!=\!8000$
releases equals the target $\rho_{\mathrm{user}}$.

\begin{table}[t]
\centering
\small
\caption{Q1: vessel Dice at matched user-level $\varepsilon$, on DRIVE
test split, mean $\pm$ standard deviation over $5$ seeds. The per-iter
model needs $\sqrt{Q}\!\approx\!90\!\times$ more noise per release to
absorb composition over $Q\!=\!8000$ patient releases, so its
feature-matching signal collapses to noise; the student is left to
learn from task loss alone. Sample-once at the same user-level
$\varepsilon$ has order-of-magnitude lower per-release noise and
preserves the distillation signal.}
\label{tab:q1-sample-once}
\begin{tabular}{r r l r l}
\toprule
$\varepsilon_{\mathrm{user}}$
   & $\sigma_{\text{per-iter}}$ & per-iter Dice
   & $\sigma_{\text{once}}$     & sample-once Dice (ours) \\
\midrule
$2$   & $5058$  & $0.567 \pm 0.011$  & $56.6$  & $\mathbf{0.581 \pm 0.023}$ \\
$4$   & $2623$  & $0.582 \pm 0.011$  & $29.3$  & $\mathbf{0.589 \pm 0.016}$ \\
$8$   & $1397$  & $0.579 \pm 0.012$  & $15.6$  & $\mathbf{0.594 \pm 0.017}$ \\
$16$  & $773$   & $0.589 \pm 0.004$  & $8.6$   & $\mathbf{0.596 \pm 0.013}$ \\
$\infty$ & $0$  & $0.676 \pm 0.007$  & $0$     & $0.676 \pm 0.007$ \\
\bottomrule
\end{tabular}
\end{table}

\paragraph{Observations.}
The Dice gap from non-private is similar across the two release models,
not because their utilities are equivalent but because at this regime
the student's task loss against the ground-truth mask dominates the
feature-matching loss: under per-iter the cached teacher signal becomes
unusable and the student falls back to its task-only baseline. The
direction of the gap is nevertheless consistent: sample-once $>$
per-iter at $4/4$ tested $\varepsilon$, with the lift growing as
$\varepsilon$ tightens. The stronger statement is in the
\emph{accounting} of Table~\ref{tab:privacy}: prior work that
reports $\varepsilon\!=\!2$ in the per-iter regime is actually
operating at user-level $\varepsilon\!=\!812$, a $400\!\times\!$
overclaim that sample-once eliminates without sacrificing utility.

% ---------------------------------------------------------------------
\subsection{Q2 — Mechanism comparison: PATE wins, CANAL bounded by $R^{1/4}$}
\label{ssec:q2-mechanism}

Table~\ref{tab:mech} compares five mechanisms on the same sample-once
framework with paired student-init testing. The first two rows isolate
the CANAL closed-form allocation against the uniform-Gaussian baseline
at $K\!=\!1$ (single teacher). Rows three and four test the proposed
PATE aggregation at $K\!=\!5$ with both noise shapes. Row five reports
the best $K$ from the saturation sweep of \S\ref{ssec:q3-k-sweep}.

\begin{table}[t]
\centering
\small
\caption{Q2: mechanism comparison on DRIVE, $5$ seeds, sample-once
framework. Paired lifts in the right block are per-seed differences
against the K=1 uniform baseline (same student initialisation, same
noise precompute seed), with standard error of the mean and the
corresponding $t$-statistic. Non-private reference Dice: $0.676$. The
bold row is our headline result. \textit{CANAL} denotes the channel-WF
allocation of Theorem~\ref{thm:wf}; PATE denotes the multi-teacher
aggregation of Theorem~\ref{thm:pate-sens}.}
\label{tab:mech}
\begin{tabular}{l l l l c c c}
\toprule
   & \multicolumn{3}{c}{vessel Dice (mean $\pm$ std)}
   & \multicolumn{3}{c}{paired lift over K=1 uniform ($t$)} \\
\cmidrule(lr){2-4}\cmidrule(lr){5-7}
mechanism
   & $\varepsilon\!=\!2$ & $\varepsilon\!=\!8$ & $\varepsilon\!=\!16$
   & $\varepsilon\!=\!2$ & $\varepsilon\!=\!8$ & $\varepsilon\!=\!16$ \\
\midrule
K=1 uniform (baseline)
   & $0.581\!\pm\!0.023$ & $0.594\!\pm\!0.017$ & $0.596\!\pm\!0.013$
   & --- & --- & --- \\
K=1 CANAL
   & $0.584\!\pm\!0.025$ & $0.593\!\pm\!0.020$ & $0.599\!\pm\!0.019$
   & $+0.003\,(\!+\!1.8\sigma)$ & $-0.000\,(\!-\!0.2\sigma)$ & $+0.003\,(\!+\!1.0\sigma)$ \\
PATE K=5 + uniform
   & $0.593\!\pm\!0.016$ & $0.618\!\pm\!0.010$ & $0.627\!\pm\!0.007$
   & $+0.012\,(\!+\!2.9\sigma)$ & $+0.025\,(\!+\!5.2\sigma)$ & $+0.031\,(\!+\!7.5\sigma)$ \\
PATE K=5 + CANAL
   & $0.592\!\pm\!0.016$ & $0.619\!\pm\!0.007$ & $0.625\!\pm\!0.007$
   & $+0.011\,(\!+\!3.0\sigma)$ & $+0.025\,(\!+\!4.2\sigma)$ & $+0.030\,(\!+\!6.3\sigma)$ \\
\midrule
\textbf{PATE K=10 + uniform}
   & $\mathbf{0.607\!\pm\!0.015}$ & $\mathbf{0.628\!\pm\!0.008}$ & $\mathbf{0.648\!\pm\!0.009}$
   & $\mathbf{+0.026\,(\!+\!5.3\sigma)}$
   & $\mathbf{+0.034\,(\!+\!8.0\sigma)}$
   & $\mathbf{+0.052\,(\!+\!15.3\sigma)}$ \\
\bottomrule
\end{tabular}
\end{table}

\paragraph{Three observations.}
\emph{(i)} \textbf{CANAL alone is empirically null on RGB fundus.} The
paired CANAL$-$uniform Dice differences at $K\!=\!1$ are within the
$\pm 0.003$ seed-noise band at every $\varepsilon$ ($|t|\!<\!2\sigma$
all three columns). This is the empirical manifestation of
Corollary~\ref{cor:r14}: with the importance ratio
$R\!\approx\!2$ on bottleneck-feature gradients, the $\sigma$ ratio
between most and least important channels is $R^{1/4}\!\approx\!1.19$,
below the Dice sensitivity of a trained decoder. CANAL remains the
optimal allocation of Theorem~\ref{thm:wf} in the importance-weighted
distortion sense; what Table~\ref{tab:mech} establishes is that its
empirical room on RGB fundus is exhausted by the $R^{1/4}$ bound.

\emph{(ii)} \textbf{PATE delivers significant Dice lift, monotonic in
$K$.} At $K\!=\!5$, PATE+uniform improves over the $K\!=\!1$ baseline by
$+0.012$ to $+0.031$ Dice ($2.9\sigma$--$7.5\sigma$) across
$\varepsilon\!\in\!\{2,8,16\}$. The headline is the bold row: at
$K\!=\!10$, the lift grows to $+0.026$ -- $+0.052$
($5.3\sigma$--$15.3\sigma$), and at $\varepsilon\!=\!16$ the distilled
student's Dice ($0.648$) \emph{exceeds} the single-teacher non-private
upper bound on the local TinyUNet ($0.620$ teacher-clean), an
ensemble-regularisation effect that is consistent with the
diversification of $K\!=\!10$ disjoint cohorts.

\emph{(iii)} \textbf{CANAL composes with PATE but does not lift it on
DRIVE.} The PATE K=5+CANAL row gives Dice within $\pm 0.0015$ of the
PATE K=5+uniform row at every $\varepsilon$, with $|t|\!<\!1.5\sigma$ on
the paired test. This is exactly the prediction of
Corollary~\ref{cor:r14}: the $R^{1/4}$ bound is independent of $K$, so
sensitivity reduction by PATE cannot create room for CANAL allocation
on a dataset where the importance spectrum is already flat. We expect
this composition to surface on data where $R$ is one to two orders of
magnitude larger; multi-modal MRI is the natural next test.

% ---------------------------------------------------------------------
\subsection{Q3 — $K$-saturation: lift grows monotonically through $K\!=\!10$}
\label{ssec:q3-k-sweep}

Theorem~\ref{thm:pate-sens} predicts $\sigma\!\propto\!1/K$. We sweep
$K\!\in\!\{1, 3, 5, 8, 10\}$ at three $\varepsilon$ values with
$5$ seeds per cell (Figure~\ref{fig:k-saturation},
Table~\ref{tab:k-sat}). Lift grows monotonically in $K$ at every
$\varepsilon$ and shows no observed saturation through $K\!=\!10$.

\begin{table}[t]
\centering
\small
\caption{Q3: PATE multi-teacher $K$-sweep. Paired lift over $K\!=\!1$
baseline; mean $\pm$ SEM of $5$ seeds, with $t$-statistic. All non-trivial
rows are highly significant. $\sigma_{\mathrm{uniform}}$ at
$\varepsilon\!=\!16$ shows the linear-in-$1/K$ scaling.}
\label{tab:k-sat}
\begin{tabular}{r r r l l l}
\toprule
   $K$ & $\Delta\!=\!2/K$
       & $\sigma$ at $\varepsilon\!=\!16$
       & lift at $\varepsilon\!=\!2$
       & lift at $\varepsilon\!=\!8$
       & lift at $\varepsilon\!=\!16$ \\
\midrule
$1$   & $2.000$ & $8.64$  & baseline $0.581$  & baseline $0.594$  & baseline $0.596$ \\
$3$   & $0.667$ & $2.88$  & $+0.008\!\pm\!0.004\,(1.7\sigma)$  & $+0.016\!\pm\!0.004\,(3.6\sigma)$  & $+0.020\!\pm\!0.003\,(6.0\sigma)$ \\
$5$   & $0.400$ & $1.73$  & $+0.012\!\pm\!0.004\,(2.9\sigma)$  & $+0.025\!\pm\!0.005\,(5.2\sigma)$  & $+0.031\!\pm\!0.004\,(7.5\sigma)$ \\
$8$   & $0.250$ & $1.08$  & $+0.020\!\pm\!0.004\,(5.0\sigma)$  & $+0.031\!\pm\!0.007\,(4.7\sigma)$  & $+0.042\!\pm\!0.005\,(8.1\sigma)$ \\
$10$  & $0.200$ & $0.86$  & $\mathbf{+0.026\!\pm\!0.005\,(5.3\sigma)}$
                        & $\mathbf{+0.034\!\pm\!0.004\,(8.0\sigma)}$
                        & $\mathbf{+0.052\!\pm\!0.003\,(15.3\sigma)}$ \\
\bottomrule
\end{tabular}
\end{table}

\begin{figure}[t]
\centering
[\textit{K-saturation figure placeholder}]
\caption{Q3: paired vessel-Dice lift over the $K\!=\!1$ baseline as a
function of the PATE ensemble size $K$ (5 seeds per cell). The Dice
lift is monotonic in $K$ at all three $\varepsilon$ values and shows
no saturation through $K\!=\!10$. The vertical separation between the
three curves reflects the better resilience of larger $\varepsilon$ to
PATE-induced noise reduction (smaller $\sigma$ matters less when
$\sigma$ is already small). Source:
\texttt{drive\_pate\_K\_saturation.png}.}
\label{fig:k-saturation}
\end{figure}

\paragraph{Cohort-size caveat.}
At $K\!=\!10$ each teacher trains on $2$ DRIVE patients. The
monotonicity through $K\!=\!10$ suggests the ensemble effect dominates
per-teacher under-training in this regime, but extrapolation to
$K\!>\!10$ would require more patients than DRIVE provides. In a
deployed federated setting each ``teacher'' is one hospital with
hundreds of patients; we report the DRIVE numbers as a controlled-
setting validation, with the understanding that the lift in
Table~\ref{tab:k-sat} is a \emph{lower bound} on what a real multi-
hospital deployment would achieve.

% ---------------------------------------------------------------------
\subsection{Limitations and scope}
\label{ssec:limitations}

\paragraph{Single dataset.} All numbers are reported on DRIVE; the
$R^{1/4}$ analysis of Corollary~\ref{cor:r14} predicts that CANAL's
empirical advantage surfaces on data with a sharper bottleneck-feature
importance spectrum (multi-modal MRI), and a feature-PATE story
generalises straightforwardly to any teacher--student feature-
distillation setting. We leave multi-dataset validation to future work.

\paragraph{Empirical privacy evaluation.} We report formal user-level
$\varepsilon$ via the zCDP accountant of Section~\ref{ssec:accounting}.
A complementary empirical evaluation by membership-inference and
feature-inversion attacks (\textit{e.g.,} LiRA \cite{carlini2022lira},
inversion \cite{dosovitskiy2016inverting}) would tighten the privacy
claim from a worst-case bound to an attack-resistant operating point.
We treat this evaluation as the principal future direction for a
follow-up.

\paragraph{Spatial allocation.} Theorem~\ref{thm:wf}'s closed form is
agnostic to the allocation index set: replacing the channel index by a
joint (channel, spatial) index leaves the formula unchanged. A
spatial-axis variant requires a data-independent saliency map, which
we sketch as a parallel future direction (\textit{e.g.,} a public-proxy
Frangi vesselness map on HRF retinal images). The local proof of
concept under a synthetic radial-bias saliency was negative, but is not
informative about the real-saliency setting; we report it only for
methodological transparency in the released code.

% ---------------------------------------------------------------------
%  BibTeX entries to add to references.bib (paste into your .bib):
%
%  @inproceedings{carlini2022lira,
%    title={Membership Inference Attacks From First Principles},
%    author={Carlini, Nicholas and others},
%    booktitle={IEEE Symposium on Security and Privacy (S\&P)},
%    year={2022},
%  }
%
%  @inproceedings{dosovitskiy2016inverting,
%    title={Inverting Visual Representations with Convolutional Networks},
%    author={Dosovitskiy, Alexey and Brox, Thomas},
%    booktitle={IEEE Conference on Computer Vision and Pattern Recognition (CVPR)},
%    year={2016},
%  }
%
%  @inproceedings{staal2004drive,
%    title={Ridge-Based Vessel Segmentation in Color Images of the Retina},
%    author={Staal, Joes and Abr\`amoff, Michael D. and Niemeijer, Meindert
%            and Viergever, Max A. and van Ginneken, Bram},
%    journal={IEEE Transactions on Medical Imaging},
%    year={2004},
%  }
%
%  @inproceedings{ronneberger2015unet,
%    title={U-Net: Convolutional Networks for Biomedical Image
%           Segmentation},
%    author={Ronneberger, Olaf and Fischer, Philipp and Brox, Thomas},
%    booktitle={Medical Image Computing and Computer-Assisted
%               Intervention (MICCAI)},
%    year={2015},
%  }
% ---------------------------------------------------------------------
 after Section 3.
%
%  Numbers in Tables 3 / 4 / 5 and Figure \ref{fig:k-saturation}
%  come from local 5-seed paired experiments in phase1_distillation/:
%
%    drive_q1_per_iter_vs_sample_once_results.json   (Tab. 3 = Q1)
%    tab4_compiled.json                              (Tab. 4 = Q2)
%    drive_pate_K_saturation_results.json            (Tab. 5 + Fig = Q3)
%    drive_pate_canal_combined_results.json          (Tab. 4 PATE+CANAL row)
%
%  All numbers in the tables can be re-extracted by running
%      python3 phase1_distillation/compile_tab4.py
% =====================================================================

\section{Experiments}
\label{sec:experiments}

We evaluate the framework of Section~\ref{sec:privacy} on retinal vessel
segmentation. Three questions structure the section:

\begin{description}
\item[(Q1)] At a \emph{matched user-level} $\varepsilon$, does the
sample-once release model preserve utility relative to the per-iteration
release model that prior work implicitly uses
(\S\ref{ssec:q1-sample-once})?
\item[(Q2)] Does the channel-importance allocation (CANAL) of
Section~\ref{ssec:wf} improve Dice over uniform noise on real DRIVE
features, and does the honest multi-teacher PATE aggregation of
Section~\ref{ssec:pate} compound with it
(\S\ref{ssec:q2-mechanism})?
\item[(Q3)] How does utility scale with the number of teachers $K$, and
where does it saturate (\S\ref{ssec:q3-k-sweep})?
\end{description}

We additionally include a paired-comparison ablation that tests
Corollary~\ref{cor:r14}'s prediction that CANAL's empirical gain on top
of PATE is bounded by $R^{1/4}$ and therefore negligible on RGB fundus
(\S\ref{ssec:q2-mechanism}, last row of Table~\ref{tab:mech}).

% ---------------------------------------------------------------------
\subsection{Experimental setup}
\label{ssec:setup}

\paragraph{Dataset.} We use DRIVE~\cite{staal2004drive}, the
canonical retinal vessel segmentation benchmark: $40$ RGB fundus images
of size $584\!\times\!565$ with expert vessel masks, split into the
canonical $20$ training and $20$ test images. Each training image is one
``user'' for the purposes of user-level DP, consistent with the
sample-once-per-image model of \S\ref{ssec:threat}.

\paragraph{Teacher and student.} Teacher and student are 3-stage
U-Net~\cite{ronneberger2015unet} variants at $96\!\times\!96$ input
resolution. The teacher has base width 32 (bottleneck dimension
$C\!=\!128$); the student has base width 16 (bottleneck dimension
$64$), with a $1\!\times\!1$ adapter projecting student features into
the teacher channel space for the feature-matching loss. Both models
are trained for $60$ epochs with Adam (lr $=\!10^{-3}$). The student
training loop runs for $40$ epochs against the sample-once cached
targets, using a combined CE+Dice task loss against the ground-truth
mask plus an MSE feature-matching loss with weight $\lambda\!=\!0.4$.

\paragraph{Privacy configuration.} We work in zCDP~\cite{bun2016concentrated}
with $\delta\!=\!10^{-5}$ throughout and report user-level
$\varepsilon\!\in\!\{2, 4, 8, 16\}$ via Eq.~\eqref{eq:eps-to-rho}.
Sensitivities follow the per-channel $L_2$ normalisation of
\S\ref{ssec:accounting}: $\Delta_c\!=\!2/K$, with $K\!=\!1$ unless
otherwise noted. The per-iteration model is parameterised by
$T\!=\!40000$, $B\!=\!4$, $N\!=\!20$, giving $Q\!=\!TB/N\!=\!8000$
per-patient releases under that threat model (Table~\ref{tab:privacy}).

\paragraph{Reporting.} Each cell of every table is the mean and standard
deviation of best validation vessel-Dice over $5$ student-init seeds
$\{100, 200, 300, 400, 500\}$. Cross-mechanism lifts are reported as
\emph{paired} per-seed differences (each seed pair shares student
initialisation, dataloader ordering, and the noise sample produced by
the sample-once precompute) with their standard error of the mean. We
report $t$-statistics on the paired lift; a $|t|\!\geq\!2$ corresponds
to $p\!<\!0.05$ in a paired one-sample test.

% ---------------------------------------------------------------------
\subsection{Q1 — sample-once preserves utility at matched user-level $\varepsilon$}
\label{ssec:q1-sample-once}

Table~\ref{tab:q1-sample-once} reports best validation vessel-Dice under
the two release models at every user-level $\varepsilon\!\in\!\{2,4,8,16\}$
that we test. The per-iter row uses the same student training loop as
sample-once, but at every training iteration adds a \emph{fresh} draw of
Gaussian noise on top of the clean teacher bottleneck of the current
batch, scaled so that the cumulative zCDP cost over $Q\!=\!8000$
releases equals the target $\rho_{\mathrm{user}}$.

\begin{table}[t]
\centering
\small
\caption{Q1: vessel Dice at matched user-level $\varepsilon$, on DRIVE
test split, mean $\pm$ standard deviation over $5$ seeds. The per-iter
model needs $\sqrt{Q}\!\approx\!90\!\times$ more noise per release to
absorb composition over $Q\!=\!8000$ patient releases, so its
feature-matching signal collapses to noise; the student is left to
learn from task loss alone. Sample-once at the same user-level
$\varepsilon$ has order-of-magnitude lower per-release noise and
preserves the distillation signal.}
\label{tab:q1-sample-once}
\begin{tabular}{r r l r l}
\toprule
$\varepsilon_{\mathrm{user}}$
   & $\sigma_{\text{per-iter}}$ & per-iter Dice
   & $\sigma_{\text{once}}$     & sample-once Dice (ours) \\
\midrule
$2$   & $5058$  & $0.567 \pm 0.011$  & $56.6$  & $\mathbf{0.581 \pm 0.023}$ \\
$4$   & $2623$  & $0.582 \pm 0.011$  & $29.3$  & $\mathbf{0.589 \pm 0.016}$ \\
$8$   & $1397$  & $0.579 \pm 0.012$  & $15.6$  & $\mathbf{0.594 \pm 0.017}$ \\
$16$  & $773$   & $0.589 \pm 0.004$  & $8.6$   & $\mathbf{0.596 \pm 0.013}$ \\
$\infty$ & $0$  & $0.676 \pm 0.007$  & $0$     & $0.676 \pm 0.007$ \\
\bottomrule
\end{tabular}
\end{table}

\paragraph{Observations.}
The Dice gap from non-private is similar across the two release models,
not because their utilities are equivalent but because at this regime
the student's task loss against the ground-truth mask dominates the
feature-matching loss: under per-iter the cached teacher signal becomes
unusable and the student falls back to its task-only baseline. The
direction of the gap is nevertheless consistent: sample-once $>$
per-iter at $4/4$ tested $\varepsilon$, with the lift growing as
$\varepsilon$ tightens. The stronger statement is in the
\emph{accounting} of Table~\ref{tab:privacy}: prior work that
reports $\varepsilon\!=\!2$ in the per-iter regime is actually
operating at user-level $\varepsilon\!=\!812$, a $400\!\times\!$
overclaim that sample-once eliminates without sacrificing utility.

% ---------------------------------------------------------------------
\subsection{Q2 — Mechanism comparison: PATE wins, CANAL bounded by $R^{1/4}$}
\label{ssec:q2-mechanism}

Table~\ref{tab:mech} compares five mechanisms on the same sample-once
framework with paired student-init testing. The first two rows isolate
the CANAL closed-form allocation against the uniform-Gaussian baseline
at $K\!=\!1$ (single teacher). Rows three and four test the proposed
PATE aggregation at $K\!=\!5$ with both noise shapes. Row five reports
the best $K$ from the saturation sweep of \S\ref{ssec:q3-k-sweep}.

\begin{table}[t]
\centering
\small
\caption{Q2: mechanism comparison on DRIVE, $5$ seeds, sample-once
framework. Paired lifts in the right block are per-seed differences
against the K=1 uniform baseline (same student initialisation, same
noise precompute seed), with standard error of the mean and the
corresponding $t$-statistic. Non-private reference Dice: $0.676$. The
bold row is our headline result. \textit{CANAL} denotes the channel-WF
allocation of Theorem~\ref{thm:wf}; PATE denotes the multi-teacher
aggregation of Theorem~\ref{thm:pate-sens}.}
\label{tab:mech}
\begin{tabular}{l l l l c c c}
\toprule
   & \multicolumn{3}{c}{vessel Dice (mean $\pm$ std)}
   & \multicolumn{3}{c}{paired lift over K=1 uniform ($t$)} \\
\cmidrule(lr){2-4}\cmidrule(lr){5-7}
mechanism
   & $\varepsilon\!=\!2$ & $\varepsilon\!=\!8$ & $\varepsilon\!=\!16$
   & $\varepsilon\!=\!2$ & $\varepsilon\!=\!8$ & $\varepsilon\!=\!16$ \\
\midrule
K=1 uniform (baseline)
   & $0.581\!\pm\!0.023$ & $0.594\!\pm\!0.017$ & $0.596\!\pm\!0.013$
   & --- & --- & --- \\
K=1 CANAL
   & $0.584\!\pm\!0.025$ & $0.593\!\pm\!0.020$ & $0.599\!\pm\!0.019$
   & $+0.003\,(\!+\!1.8\sigma)$ & $-0.000\,(\!-\!0.2\sigma)$ & $+0.003\,(\!+\!1.0\sigma)$ \\
PATE K=5 + uniform
   & $0.593\!\pm\!0.016$ & $0.618\!\pm\!0.010$ & $0.627\!\pm\!0.007$
   & $+0.012\,(\!+\!2.9\sigma)$ & $+0.025\,(\!+\!5.2\sigma)$ & $+0.031\,(\!+\!7.5\sigma)$ \\
PATE K=5 + CANAL
   & $0.592\!\pm\!0.016$ & $0.619\!\pm\!0.007$ & $0.625\!\pm\!0.007$
   & $+0.011\,(\!+\!3.0\sigma)$ & $+0.025\,(\!+\!4.2\sigma)$ & $+0.030\,(\!+\!6.3\sigma)$ \\
\midrule
\textbf{PATE K=10 + uniform}
   & $\mathbf{0.607\!\pm\!0.015}$ & $\mathbf{0.628\!\pm\!0.008}$ & $\mathbf{0.648\!\pm\!0.009}$
   & $\mathbf{+0.026\,(\!+\!5.3\sigma)}$
   & $\mathbf{+0.034\,(\!+\!8.0\sigma)}$
   & $\mathbf{+0.052\,(\!+\!15.3\sigma)}$ \\
\bottomrule
\end{tabular}
\end{table}

\paragraph{Three observations.}
\emph{(i)} \textbf{CANAL alone is empirically null on RGB fundus.} The
paired CANAL$-$uniform Dice differences at $K\!=\!1$ are within the
$\pm 0.003$ seed-noise band at every $\varepsilon$ ($|t|\!<\!2\sigma$
all three columns). This is the empirical manifestation of
Corollary~\ref{cor:r14}: with the importance ratio
$R\!\approx\!2$ on bottleneck-feature gradients, the $\sigma$ ratio
between most and least important channels is $R^{1/4}\!\approx\!1.19$,
below the Dice sensitivity of a trained decoder. CANAL remains the
optimal allocation of Theorem~\ref{thm:wf} in the importance-weighted
distortion sense; what Table~\ref{tab:mech} establishes is that its
empirical room on RGB fundus is exhausted by the $R^{1/4}$ bound.

\emph{(ii)} \textbf{PATE delivers significant Dice lift, monotonic in
$K$.} At $K\!=\!5$, PATE+uniform improves over the $K\!=\!1$ baseline by
$+0.012$ to $+0.031$ Dice ($2.9\sigma$--$7.5\sigma$) across
$\varepsilon\!\in\!\{2,8,16\}$. The headline is the bold row: at
$K\!=\!10$, the lift grows to $+0.026$ -- $+0.052$
($5.3\sigma$--$15.3\sigma$), and at $\varepsilon\!=\!16$ the distilled
student's Dice ($0.648$) \emph{exceeds} the single-teacher non-private
upper bound on the local TinyUNet ($0.620$ teacher-clean), an
ensemble-regularisation effect that is consistent with the
diversification of $K\!=\!10$ disjoint cohorts.

\emph{(iii)} \textbf{CANAL composes with PATE but does not lift it on
DRIVE.} The PATE K=5+CANAL row gives Dice within $\pm 0.0015$ of the
PATE K=5+uniform row at every $\varepsilon$, with $|t|\!<\!1.5\sigma$ on
the paired test. This is exactly the prediction of
Corollary~\ref{cor:r14}: the $R^{1/4}$ bound is independent of $K$, so
sensitivity reduction by PATE cannot create room for CANAL allocation
on a dataset where the importance spectrum is already flat. We expect
this composition to surface on data where $R$ is one to two orders of
magnitude larger; multi-modal MRI is the natural next test.

% ---------------------------------------------------------------------
\subsection{Q3 — $K$-saturation: lift grows monotonically through $K\!=\!10$}
\label{ssec:q3-k-sweep}

Theorem~\ref{thm:pate-sens} predicts $\sigma\!\propto\!1/K$. We sweep
$K\!\in\!\{1, 3, 5, 8, 10\}$ at three $\varepsilon$ values with
$5$ seeds per cell (Figure~\ref{fig:k-saturation},
Table~\ref{tab:k-sat}). Lift grows monotonically in $K$ at every
$\varepsilon$ and shows no observed saturation through $K\!=\!10$.

\begin{table}[t]
\centering
\small
\caption{Q3: PATE multi-teacher $K$-sweep. Paired lift over $K\!=\!1$
baseline; mean $\pm$ SEM of $5$ seeds, with $t$-statistic. All non-trivial
rows are highly significant. $\sigma_{\mathrm{uniform}}$ at
$\varepsilon\!=\!16$ shows the linear-in-$1/K$ scaling.}
\label{tab:k-sat}
\begin{tabular}{r r r l l l}
\toprule
   $K$ & $\Delta\!=\!2/K$
       & $\sigma$ at $\varepsilon\!=\!16$
       & lift at $\varepsilon\!=\!2$
       & lift at $\varepsilon\!=\!8$
       & lift at $\varepsilon\!=\!16$ \\
\midrule
$1$   & $2.000$ & $8.64$  & baseline $0.581$  & baseline $0.594$  & baseline $0.596$ \\
$3$   & $0.667$ & $2.88$  & $+0.008\!\pm\!0.004\,(1.7\sigma)$  & $+0.016\!\pm\!0.004\,(3.6\sigma)$  & $+0.020\!\pm\!0.003\,(6.0\sigma)$ \\
$5$   & $0.400$ & $1.73$  & $+0.012\!\pm\!0.004\,(2.9\sigma)$  & $+0.025\!\pm\!0.005\,(5.2\sigma)$  & $+0.031\!\pm\!0.004\,(7.5\sigma)$ \\
$8$   & $0.250$ & $1.08$  & $+0.020\!\pm\!0.004\,(5.0\sigma)$  & $+0.031\!\pm\!0.007\,(4.7\sigma)$  & $+0.042\!\pm\!0.005\,(8.1\sigma)$ \\
$10$  & $0.200$ & $0.86$  & $\mathbf{+0.026\!\pm\!0.005\,(5.3\sigma)}$
                        & $\mathbf{+0.034\!\pm\!0.004\,(8.0\sigma)}$
                        & $\mathbf{+0.052\!\pm\!0.003\,(15.3\sigma)}$ \\
\bottomrule
\end{tabular}
\end{table}

\begin{figure}[t]
\centering
[\textit{K-saturation figure placeholder}]
\caption{Q3: paired vessel-Dice lift over the $K\!=\!1$ baseline as a
function of the PATE ensemble size $K$ (5 seeds per cell). The Dice
lift is monotonic in $K$ at all three $\varepsilon$ values and shows
no saturation through $K\!=\!10$. The vertical separation between the
three curves reflects the better resilience of larger $\varepsilon$ to
PATE-induced noise reduction (smaller $\sigma$ matters less when
$\sigma$ is already small). Source:
\texttt{drive\_pate\_K\_saturation.png}.}
\label{fig:k-saturation}
\end{figure}

\paragraph{Cohort-size caveat.}
At $K\!=\!10$ each teacher trains on $2$ DRIVE patients. The
monotonicity through $K\!=\!10$ suggests the ensemble effect dominates
per-teacher under-training in this regime, but extrapolation to
$K\!>\!10$ would require more patients than DRIVE provides. In a
deployed federated setting each ``teacher'' is one hospital with
hundreds of patients; we report the DRIVE numbers as a controlled-
setting validation, with the understanding that the lift in
Table~\ref{tab:k-sat} is a \emph{lower bound} on what a real multi-
hospital deployment would achieve.

% ---------------------------------------------------------------------
\subsection{Limitations and scope}
\label{ssec:limitations}

\paragraph{Single dataset.} All numbers are reported on DRIVE; the
$R^{1/4}$ analysis of Corollary~\ref{cor:r14} predicts that CANAL's
empirical advantage surfaces on data with a sharper bottleneck-feature
importance spectrum (multi-modal MRI), and a feature-PATE story
generalises straightforwardly to any teacher--student feature-
distillation setting. We leave multi-dataset validation to future work.

\paragraph{Empirical privacy evaluation.} We report formal user-level
$\varepsilon$ via the zCDP accountant of Section~\ref{ssec:accounting}.
A complementary empirical evaluation by membership-inference and
feature-inversion attacks (\textit{e.g.,} LiRA \cite{carlini2022lira},
inversion \cite{dosovitskiy2016inverting}) would tighten the privacy
claim from a worst-case bound to an attack-resistant operating point.
We treat this evaluation as the principal future direction for a
follow-up.

\paragraph{Spatial allocation.} Theorem~\ref{thm:wf}'s closed form is
agnostic to the allocation index set: replacing the channel index by a
joint (channel, spatial) index leaves the formula unchanged. A
spatial-axis variant requires a data-independent saliency map, which
we sketch as a parallel future direction (\textit{e.g.,} a public-proxy
Frangi vesselness map on HRF retinal images). The local proof of
concept under a synthetic radial-bias saliency was negative, but is not
informative about the real-saliency setting; we report it only for
methodological transparency in the released code.

% ---------------------------------------------------------------------
%  BibTeX entries to add to references.bib (paste into your .bib):
%
%  @inproceedings{carlini2022lira,
%    title={Membership Inference Attacks From First Principles},
%    author={Carlini, Nicholas and others},
%    booktitle={IEEE Symposium on Security and Privacy (S\&P)},
%    year={2022},
%  }
%
%  @inproceedings{dosovitskiy2016inverting,
%    title={Inverting Visual Representations with Convolutional Networks},
%    author={Dosovitskiy, Alexey and Brox, Thomas},
%    booktitle={IEEE Conference on Computer Vision and Pattern Recognition (CVPR)},
%    year={2016},
%  }
%
%  @inproceedings{staal2004drive,
%    title={Ridge-Based Vessel Segmentation in Color Images of the Retina},
%    author={Staal, Joes and Abr\`amoff, Michael D. and Niemeijer, Meindert
%            and Viergever, Max A. and van Ginneken, Bram},
%    journal={IEEE Transactions on Medical Imaging},
%    year={2004},
%  }
%
%  @inproceedings{ronneberger2015unet,
%    title={U-Net: Convolutional Networks for Biomedical Image
%           Segmentation},
%    author={Ronneberger, Olaf and Fischer, Philipp and Brox, Thomas},
%    booktitle={Medical Image Computing and Computer-Assisted
%               Intervention (MICCAI)},
%    year={2015},
%  }
% ---------------------------------------------------------------------
 after Section 3.
%
%  Numbers in Tables 3 / 4 / 5 and Figure \ref{fig:k-saturation}
%  come from local 5-seed paired experiments in phase1_distillation/:
%
%    drive_q1_per_iter_vs_sample_once_results.json   (Tab. 3 = Q1)
%    tab4_compiled.json                              (Tab. 4 = Q2)
%    drive_pate_K_saturation_results.json            (Tab. 5 + Fig = Q3)
%    drive_pate_canal_combined_results.json          (Tab. 4 PATE+CANAL row)
%
%  All numbers in the tables can be re-extracted by running
%      python3 phase1_distillation/compile_tab4.py
% =====================================================================

\section{Experiments}
\label{sec:experiments}

We evaluate the framework of Section~\ref{sec:privacy} on retinal vessel
segmentation. Three questions structure the section:

\begin{description}
\item[(Q1)] At a \emph{matched user-level} $\varepsilon$, does the
sample-once release model preserve utility relative to the per-iteration
release model that prior work implicitly uses
(\S\ref{ssec:q1-sample-once})?
\item[(Q2)] Does the channel-importance allocation (CANAL) of
Section~\ref{ssec:wf} improve Dice over uniform noise on real DRIVE
features, and does the honest multi-teacher PATE aggregation of
Section~\ref{ssec:pate} compound with it
(\S\ref{ssec:q2-mechanism})?
\item[(Q3)] How does utility scale with the number of teachers $K$, and
where does it saturate (\S\ref{ssec:q3-k-sweep})?
\end{description}

We additionally include a paired-comparison ablation that tests
Corollary~\ref{cor:r14}'s prediction that CANAL's empirical gain on top
of PATE is bounded by $R^{1/4}$ and therefore negligible on RGB fundus
(\S\ref{ssec:q2-mechanism}, last row of Table~\ref{tab:mech}).

% ---------------------------------------------------------------------
\subsection{Experimental setup}
\label{ssec:setup}

\paragraph{Dataset.} We use DRIVE~\cite{staal2004drive}, the
canonical retinal vessel segmentation benchmark: $40$ RGB fundus images
of size $584\!\times\!565$ with expert vessel masks, split into the
canonical $20$ training and $20$ test images. Each training image is one
``user'' for the purposes of user-level DP, consistent with the
sample-once-per-image model of \S\ref{ssec:threat}.

\paragraph{Teacher and student.} Teacher and student are 3-stage
U-Net~\cite{ronneberger2015unet} variants at $96\!\times\!96$ input
resolution. The teacher has base width 32 (bottleneck dimension
$C\!=\!128$); the student has base width 16 (bottleneck dimension
$64$), with a $1\!\times\!1$ adapter projecting student features into
the teacher channel space for the feature-matching loss. Both models
are trained for $60$ epochs with Adam (lr $=\!10^{-3}$). The student
training loop runs for $40$ epochs against the sample-once cached
targets, using a combined CE+Dice task loss against the ground-truth
mask plus an MSE feature-matching loss with weight $\lambda\!=\!0.4$.

\paragraph{Privacy configuration.} We work in zCDP~\cite{bun2016concentrated}
with $\delta\!=\!10^{-5}$ throughout and report user-level
$\varepsilon\!\in\!\{2, 4, 8, 16\}$ via Eq.~\eqref{eq:eps-to-rho}.
Sensitivities follow the per-channel $L_2$ normalisation of
\S\ref{ssec:accounting}: $\Delta_c\!=\!2/K$, with $K\!=\!1$ unless
otherwise noted. The per-iteration model is parameterised by
$T\!=\!40000$, $B\!=\!4$, $N\!=\!20$, giving $Q\!=\!TB/N\!=\!8000$
per-patient releases under that threat model (Table~\ref{tab:privacy}).

\paragraph{Reporting.} Each cell of every table is the mean and standard
deviation of best validation vessel-Dice over $5$ student-init seeds
$\{100, 200, 300, 400, 500\}$. Cross-mechanism lifts are reported as
\emph{paired} per-seed differences (each seed pair shares student
initialisation, dataloader ordering, and the noise sample produced by
the sample-once precompute) with their standard error of the mean. We
report $t$-statistics on the paired lift; a $|t|\!\geq\!2$ corresponds
to $p\!<\!0.05$ in a paired one-sample test.

% ---------------------------------------------------------------------
\subsection{Q1 — sample-once preserves utility at matched user-level $\varepsilon$}
\label{ssec:q1-sample-once}

Table~\ref{tab:q1-sample-once} reports best validation vessel-Dice under
the two release models at every user-level $\varepsilon\!\in\!\{2,4,8,16\}$
that we test. The per-iter row uses the same student training loop as
sample-once, but at every training iteration adds a \emph{fresh} draw of
Gaussian noise on top of the clean teacher bottleneck of the current
batch, scaled so that the cumulative zCDP cost over $Q\!=\!8000$
releases equals the target $\rho_{\mathrm{user}}$.

\begin{table}[t]
\centering
\small
\caption{Q1: vessel Dice at matched user-level $\varepsilon$, on DRIVE
test split, mean $\pm$ standard deviation over $5$ seeds. The per-iter
model needs $\sqrt{Q}\!\approx\!90\!\times$ more noise per release to
absorb composition over $Q\!=\!8000$ patient releases, so its
feature-matching signal collapses to noise; the student is left to
learn from task loss alone. Sample-once at the same user-level
$\varepsilon$ has order-of-magnitude lower per-release noise and
preserves the distillation signal.}
\label{tab:q1-sample-once}
\begin{tabular}{r r l r l}
\toprule
$\varepsilon_{\mathrm{user}}$
   & $\sigma_{\text{per-iter}}$ & per-iter Dice
   & $\sigma_{\text{once}}$     & sample-once Dice (ours) \\
\midrule
$2$   & $5058$  & $0.567 \pm 0.011$  & $56.6$  & $\mathbf{0.581 \pm 0.023}$ \\
$4$   & $2623$  & $0.582 \pm 0.011$  & $29.3$  & $\mathbf{0.589 \pm 0.016}$ \\
$8$   & $1397$  & $0.579 \pm 0.012$  & $15.6$  & $\mathbf{0.594 \pm 0.017}$ \\
$16$  & $773$   & $0.589 \pm 0.004$  & $8.6$   & $\mathbf{0.596 \pm 0.013}$ \\
$\infty$ & $0$  & $0.676 \pm 0.007$  & $0$     & $0.676 \pm 0.007$ \\
\bottomrule
\end{tabular}
\end{table}

\paragraph{Observations.}
The Dice gap from non-private is similar across the two release models,
not because their utilities are equivalent but because at this regime
the student's task loss against the ground-truth mask dominates the
feature-matching loss: under per-iter the cached teacher signal becomes
unusable and the student falls back to its task-only baseline. The
direction of the gap is nevertheless consistent: sample-once $>$
per-iter at $4/4$ tested $\varepsilon$, with the lift growing as
$\varepsilon$ tightens. The stronger statement is in the
\emph{accounting} of Table~\ref{tab:privacy}: prior work that
reports $\varepsilon\!=\!2$ in the per-iter regime is actually
operating at user-level $\varepsilon\!=\!812$, a $400\!\times\!$
overclaim that sample-once eliminates without sacrificing utility.

% ---------------------------------------------------------------------
\subsection{Q2 — Mechanism comparison: PATE wins, CANAL bounded by $R^{1/4}$}
\label{ssec:q2-mechanism}

Table~\ref{tab:mech} compares five mechanisms on the same sample-once
framework with paired student-init testing. The first two rows isolate
the CANAL closed-form allocation against the uniform-Gaussian baseline
at $K\!=\!1$ (single teacher). Rows three and four test the proposed
PATE aggregation at $K\!=\!5$ with both noise shapes. Row five reports
the best $K$ from the saturation sweep of \S\ref{ssec:q3-k-sweep}.

\begin{table}[t]
\centering
\small
\caption{Q2: mechanism comparison on DRIVE, $5$ seeds, sample-once
framework. Paired lifts in the right block are per-seed differences
against the K=1 uniform baseline (same student initialisation, same
noise precompute seed), with standard error of the mean and the
corresponding $t$-statistic. Non-private reference Dice: $0.676$. The
bold row is our headline result. \textit{CANAL} denotes the channel-WF
allocation of Theorem~\ref{thm:wf}; PATE denotes the multi-teacher
aggregation of Theorem~\ref{thm:pate-sens}.}
\label{tab:mech}
\begin{tabular}{l l l l c c c}
\toprule
   & \multicolumn{3}{c}{vessel Dice (mean $\pm$ std)}
   & \multicolumn{3}{c}{paired lift over K=1 uniform ($t$)} \\
\cmidrule(lr){2-4}\cmidrule(lr){5-7}
mechanism
   & $\varepsilon\!=\!2$ & $\varepsilon\!=\!8$ & $\varepsilon\!=\!16$
   & $\varepsilon\!=\!2$ & $\varepsilon\!=\!8$ & $\varepsilon\!=\!16$ \\
\midrule
K=1 uniform (baseline)
   & $0.581\!\pm\!0.023$ & $0.594\!\pm\!0.017$ & $0.596\!\pm\!0.013$
   & --- & --- & --- \\
K=1 CANAL
   & $0.584\!\pm\!0.025$ & $0.593\!\pm\!0.020$ & $0.599\!\pm\!0.019$
   & $+0.003\,(\!+\!1.8\sigma)$ & $-0.000\,(\!-\!0.2\sigma)$ & $+0.003\,(\!+\!1.0\sigma)$ \\
PATE K=5 + uniform
   & $0.593\!\pm\!0.016$ & $0.618\!\pm\!0.010$ & $0.627\!\pm\!0.007$
   & $+0.012\,(\!+\!2.9\sigma)$ & $+0.025\,(\!+\!5.2\sigma)$ & $+0.031\,(\!+\!7.5\sigma)$ \\
PATE K=5 + CANAL
   & $0.592\!\pm\!0.016$ & $0.619\!\pm\!0.007$ & $0.625\!\pm\!0.007$
   & $+0.011\,(\!+\!3.0\sigma)$ & $+0.025\,(\!+\!4.2\sigma)$ & $+0.030\,(\!+\!6.3\sigma)$ \\
\midrule
\textbf{PATE K=10 + uniform}
   & $\mathbf{0.607\!\pm\!0.015}$ & $\mathbf{0.628\!\pm\!0.008}$ & $\mathbf{0.648\!\pm\!0.009}$
   & $\mathbf{+0.026\,(\!+\!5.3\sigma)}$
   & $\mathbf{+0.034\,(\!+\!8.0\sigma)}$
   & $\mathbf{+0.052\,(\!+\!15.3\sigma)}$ \\
\bottomrule
\end{tabular}
\end{table}

\paragraph{Three observations.}
\emph{(i)} \textbf{CANAL alone is empirically null on RGB fundus.} The
paired CANAL$-$uniform Dice differences at $K\!=\!1$ are within the
$\pm 0.003$ seed-noise band at every $\varepsilon$ ($|t|\!<\!2\sigma$
all three columns). This is the empirical manifestation of
Corollary~\ref{cor:r14}: with the importance ratio
$R\!\approx\!2$ on bottleneck-feature gradients, the $\sigma$ ratio
between most and least important channels is $R^{1/4}\!\approx\!1.19$,
below the Dice sensitivity of a trained decoder. CANAL remains the
optimal allocation of Theorem~\ref{thm:wf} in the importance-weighted
distortion sense; what Table~\ref{tab:mech} establishes is that its
empirical room on RGB fundus is exhausted by the $R^{1/4}$ bound.

\emph{(ii)} \textbf{PATE delivers significant Dice lift, monotonic in
$K$.} At $K\!=\!5$, PATE+uniform improves over the $K\!=\!1$ baseline by
$+0.012$ to $+0.031$ Dice ($2.9\sigma$--$7.5\sigma$) across
$\varepsilon\!\in\!\{2,8,16\}$. The headline is the bold row: at
$K\!=\!10$, the lift grows to $+0.026$ -- $+0.052$
($5.3\sigma$--$15.3\sigma$), and at $\varepsilon\!=\!16$ the distilled
student's Dice ($0.648$) \emph{exceeds} the single-teacher non-private
upper bound on the local TinyUNet ($0.620$ teacher-clean), an
ensemble-regularisation effect that is consistent with the
diversification of $K\!=\!10$ disjoint cohorts.

\emph{(iii)} \textbf{CANAL composes with PATE but does not lift it on
DRIVE.} The PATE K=5+CANAL row gives Dice within $\pm 0.0015$ of the
PATE K=5+uniform row at every $\varepsilon$, with $|t|\!<\!1.5\sigma$ on
the paired test. This is exactly the prediction of
Corollary~\ref{cor:r14}: the $R^{1/4}$ bound is independent of $K$, so
sensitivity reduction by PATE cannot create room for CANAL allocation
on a dataset where the importance spectrum is already flat. We expect
this composition to surface on data where $R$ is one to two orders of
magnitude larger; multi-modal MRI is the natural next test.

% ---------------------------------------------------------------------
\subsection{Q3 — $K$-saturation: lift grows monotonically through $K\!=\!10$}
\label{ssec:q3-k-sweep}

Theorem~\ref{thm:pate-sens} predicts $\sigma\!\propto\!1/K$. We sweep
$K\!\in\!\{1, 3, 5, 8, 10\}$ at three $\varepsilon$ values with
$5$ seeds per cell (Figure~\ref{fig:k-saturation},
Table~\ref{tab:k-sat}). Lift grows monotonically in $K$ at every
$\varepsilon$ and shows no observed saturation through $K\!=\!10$.

\begin{table}[t]
\centering
\small
\caption{Q3: PATE multi-teacher $K$-sweep. Paired lift over $K\!=\!1$
baseline; mean $\pm$ SEM of $5$ seeds, with $t$-statistic. All non-trivial
rows are highly significant. $\sigma_{\mathrm{uniform}}$ at
$\varepsilon\!=\!16$ shows the linear-in-$1/K$ scaling.}
\label{tab:k-sat}
\begin{tabular}{r r r l l l}
\toprule
   $K$ & $\Delta\!=\!2/K$
       & $\sigma$ at $\varepsilon\!=\!16$
       & lift at $\varepsilon\!=\!2$
       & lift at $\varepsilon\!=\!8$
       & lift at $\varepsilon\!=\!16$ \\
\midrule
$1$   & $2.000$ & $8.64$  & baseline $0.581$  & baseline $0.594$  & baseline $0.596$ \\
$3$   & $0.667$ & $2.88$  & $+0.008\!\pm\!0.004\,(1.7\sigma)$  & $+0.016\!\pm\!0.004\,(3.6\sigma)$  & $+0.020\!\pm\!0.003\,(6.0\sigma)$ \\
$5$   & $0.400$ & $1.73$  & $+0.012\!\pm\!0.004\,(2.9\sigma)$  & $+0.025\!\pm\!0.005\,(5.2\sigma)$  & $+0.031\!\pm\!0.004\,(7.5\sigma)$ \\
$8$   & $0.250$ & $1.08$  & $+0.020\!\pm\!0.004\,(5.0\sigma)$  & $+0.031\!\pm\!0.007\,(4.7\sigma)$  & $+0.042\!\pm\!0.005\,(8.1\sigma)$ \\
$10$  & $0.200$ & $0.86$  & $\mathbf{+0.026\!\pm\!0.005\,(5.3\sigma)}$
                        & $\mathbf{+0.034\!\pm\!0.004\,(8.0\sigma)}$
                        & $\mathbf{+0.052\!\pm\!0.003\,(15.3\sigma)}$ \\
\bottomrule
\end{tabular}
\end{table}

\begin{figure}[t]
\centering
[\textit{K-saturation figure placeholder}]
\caption{Q3: paired vessel-Dice lift over the $K\!=\!1$ baseline as a
function of the PATE ensemble size $K$ (5 seeds per cell). The Dice
lift is monotonic in $K$ at all three $\varepsilon$ values and shows
no saturation through $K\!=\!10$. The vertical separation between the
three curves reflects the better resilience of larger $\varepsilon$ to
PATE-induced noise reduction (smaller $\sigma$ matters less when
$\sigma$ is already small). Source:
\texttt{drive\_pate\_K\_saturation.png}.}
\label{fig:k-saturation}
\end{figure}

\paragraph{Cohort-size caveat.}
At $K\!=\!10$ each teacher trains on $2$ DRIVE patients. The
monotonicity through $K\!=\!10$ suggests the ensemble effect dominates
per-teacher under-training in this regime, but extrapolation to
$K\!>\!10$ would require more patients than DRIVE provides. In a
deployed federated setting each ``teacher'' is one hospital with
hundreds of patients; we report the DRIVE numbers as a controlled-
setting validation, with the understanding that the lift in
Table~\ref{tab:k-sat} is a \emph{lower bound} on what a real multi-
hospital deployment would achieve.

% ---------------------------------------------------------------------
\subsection{Limitations and scope}
\label{ssec:limitations}

\paragraph{Single dataset.} All numbers are reported on DRIVE; the
$R^{1/4}$ analysis of Corollary~\ref{cor:r14} predicts that CANAL's
empirical advantage surfaces on data with a sharper bottleneck-feature
importance spectrum (multi-modal MRI), and a feature-PATE story
generalises straightforwardly to any teacher--student feature-
distillation setting. We leave multi-dataset validation to future work.

\paragraph{Empirical privacy evaluation.} We report formal user-level
$\varepsilon$ via the zCDP accountant of Section~\ref{ssec:accounting}.
A complementary empirical evaluation by membership-inference and
feature-inversion attacks (\textit{e.g.,} LiRA \cite{carlini2022lira},
inversion \cite{dosovitskiy2016inverting}) would tighten the privacy
claim from a worst-case bound to an attack-resistant operating point.
We treat this evaluation as the principal future direction for a
follow-up.

\paragraph{Spatial allocation.} Theorem~\ref{thm:wf}'s closed form is
agnostic to the allocation index set: replacing the channel index by a
joint (channel, spatial) index leaves the formula unchanged. A
spatial-axis variant requires a data-independent saliency map, which
we sketch as a parallel future direction (\textit{e.g.,} a public-proxy
Frangi vesselness map on HRF retinal images). The local proof of
concept under a synthetic radial-bias saliency was negative, but is not
informative about the real-saliency setting; we report it only for
methodological transparency in the released code.

% ---------------------------------------------------------------------
%  BibTeX entries to add to references.bib (paste into your .bib):
%
%  @inproceedings{carlini2022lira,
%    title={Membership Inference Attacks From First Principles},
%    author={Carlini, Nicholas and others},
%    booktitle={IEEE Symposium on Security and Privacy (S\&P)},
%    year={2022},
%  }
%
%  @inproceedings{dosovitskiy2016inverting,
%    title={Inverting Visual Representations with Convolutional Networks},
%    author={Dosovitskiy, Alexey and Brox, Thomas},
%    booktitle={IEEE Conference on Computer Vision and Pattern Recognition (CVPR)},
%    year={2016},
%  }
%
%  @inproceedings{staal2004drive,
%    title={Ridge-Based Vessel Segmentation in Color Images of the Retina},
%    author={Staal, Joes and Abr\`amoff, Michael D. and Niemeijer, Meindert
%            and Viergever, Max A. and van Ginneken, Bram},
%    journal={IEEE Transactions on Medical Imaging},
%    year={2004},
%  }
%
%  @inproceedings{ronneberger2015unet,
%    title={U-Net: Convolutional Networks for Biomedical Image
%           Segmentation},
%    author={Ronneberger, Olaf and Fischer, Philipp and Brox, Thomas},
%    booktitle={Medical Image Computing and Computer-Assisted
%               Intervention (MICCAI)},
%    year={2015},
%  }
% ---------------------------------------------------------------------
 after Section 3.
%
%  Numbers in Tables 3 / 4 / 5 and Figure \ref{fig:k-saturation}
%  come from local 5-seed paired experiments in phase1_distillation/:
%
%    drive_q1_per_iter_vs_sample_once_results.json   (Tab. 3 = Q1)
%    tab4_compiled.json                              (Tab. 4 = Q2)
%    drive_pate_K_saturation_results.json            (Tab. 5 + Fig = Q3)
%    drive_pate_canal_combined_results.json          (Tab. 4 PATE+CANAL row)
%
%  All numbers in the tables can be re-extracted by running
%      python3 phase1_distillation/compile_tab4.py
% =====================================================================

\section{Experiments}
\label{sec:experiments}

We evaluate the framework of Section~\ref{sec:privacy} on retinal vessel
segmentation. Three questions structure the section:

\begin{description}
\item[(Q1)] At a \emph{matched user-level} $\varepsilon$, does the
sample-once release model preserve utility relative to the per-iteration
release model that prior work implicitly uses
(\S\ref{ssec:q1-sample-once})?
\item[(Q2)] Does the channel-importance allocation (CANAL) of
Section~\ref{ssec:wf} improve Dice over uniform noise on real DRIVE
features, and does the honest multi-teacher PATE aggregation of
Section~\ref{ssec:pate} compound with it
(\S\ref{ssec:q2-mechanism})?
\item[(Q3)] How does utility scale with the number of teachers $K$, and
where does it saturate (\S\ref{ssec:q3-k-sweep})?
\end{description}

We additionally include a paired-comparison ablation that tests
Corollary~\ref{cor:r14}'s prediction that CANAL's empirical gain on top
of PATE is bounded by $R^{1/4}$ and therefore negligible on RGB fundus
(\S\ref{ssec:q2-mechanism}, last row of Table~\ref{tab:mech}).

% ---------------------------------------------------------------------
\subsection{Experimental setup}
\label{ssec:setup}

\paragraph{Dataset.} We use DRIVE~\cite{staal2004drive}, the
canonical retinal vessel segmentation benchmark: $40$ RGB fundus images
of size $584\!\times\!565$ with expert vessel masks, split into the
canonical $20$ training and $20$ test images. Each training image is one
``user'' for the purposes of user-level DP, consistent with the
sample-once-per-image model of \S\ref{ssec:threat}.

\paragraph{Teacher and student.} Teacher and student are 3-stage
U-Net~\cite{ronneberger2015unet} variants at $96\!\times\!96$ input
resolution. The teacher has base width 32 (bottleneck dimension
$C\!=\!128$); the student has base width 16 (bottleneck dimension
$64$), with a $1\!\times\!1$ adapter projecting student features into
the teacher channel space for the feature-matching loss. Both models
are trained for $60$ epochs with Adam (lr $=\!10^{-3}$). The student
training loop runs for $40$ epochs against the sample-once cached
targets, using a combined CE+Dice task loss against the ground-truth
mask plus an MSE feature-matching loss with weight $\lambda\!=\!0.4$.

\paragraph{Privacy configuration.} We work in zCDP~\cite{bun2016concentrated}
with $\delta\!=\!10^{-5}$ throughout and report user-level
$\varepsilon\!\in\!\{2, 4, 8, 16\}$ via Eq.~\eqref{eq:eps-to-rho}.
Sensitivities follow the per-channel $L_2$ normalisation of
\S\ref{ssec:accounting}: $\Delta_c\!=\!2/K$, with $K\!=\!1$ unless
otherwise noted. The per-iteration model is parameterised by
$T\!=\!40000$, $B\!=\!4$, $N\!=\!20$, giving $Q\!=\!TB/N\!=\!8000$
per-patient releases under that threat model (Table~\ref{tab:privacy}).

\paragraph{Reporting.} Each cell of every table is the mean and standard
deviation of best validation vessel-Dice over $5$ student-init seeds
$\{100, 200, 300, 400, 500\}$. Cross-mechanism lifts are reported as
\emph{paired} per-seed differences (each seed pair shares student
initialisation, dataloader ordering, and the noise sample produced by
the sample-once precompute) with their standard error of the mean. We
report $t$-statistics on the paired lift; a $|t|\!\geq\!2$ corresponds
to $p\!<\!0.05$ in a paired one-sample test.

% ---------------------------------------------------------------------
\subsection{Q1 — sample-once preserves utility at matched user-level $\varepsilon$}
\label{ssec:q1-sample-once}

Table~\ref{tab:q1-sample-once} reports best validation vessel-Dice under
the two release models at every user-level $\varepsilon\!\in\!\{2,4,8,16\}$
that we test. The per-iter row uses the same student training loop as
sample-once, but at every training iteration adds a \emph{fresh} draw of
Gaussian noise on top of the clean teacher bottleneck of the current
batch, scaled so that the cumulative zCDP cost over $Q\!=\!8000$
releases equals the target $\rho_{\mathrm{user}}$.

\begin{table}[t]
\centering
\small
\caption{Q1: vessel Dice at matched user-level $\varepsilon$, on DRIVE
test split, mean $\pm$ standard deviation over $5$ seeds. The per-iter
model needs $\sqrt{Q}\!\approx\!90\!\times$ more noise per release to
absorb composition over $Q\!=\!8000$ patient releases, so its
feature-matching signal collapses to noise; the student is left to
learn from task loss alone. Sample-once at the same user-level
$\varepsilon$ has order-of-magnitude lower per-release noise and
preserves the distillation signal.}
\label{tab:q1-sample-once}
\begin{tabular}{r r l r l}
\toprule
$\varepsilon_{\mathrm{user}}$
   & $\sigma_{\text{per-iter}}$ & per-iter Dice
   & $\sigma_{\text{once}}$     & sample-once Dice (ours) \\
\midrule
$2$   & $5058$  & $0.567 \pm 0.011$  & $56.6$  & $\mathbf{0.581 \pm 0.023}$ \\
$4$   & $2623$  & $0.582 \pm 0.011$  & $29.3$  & $\mathbf{0.589 \pm 0.016}$ \\
$8$   & $1397$  & $0.579 \pm 0.012$  & $15.6$  & $\mathbf{0.594 \pm 0.017}$ \\
$16$  & $773$   & $0.589 \pm 0.004$  & $8.6$   & $\mathbf{0.596 \pm 0.013}$ \\
$\infty$ & $0$  & $0.676 \pm 0.007$  & $0$     & $0.676 \pm 0.007$ \\
\bottomrule
\end{tabular}
\end{table}

\paragraph{Observations.}
The Dice gap from non-private is similar across the two release models,
not because their utilities are equivalent but because at this regime
the student's task loss against the ground-truth mask dominates the
feature-matching loss: under per-iter the cached teacher signal becomes
unusable and the student falls back to its task-only baseline. The
direction of the gap is nevertheless consistent: sample-once $>$
per-iter at $4/4$ tested $\varepsilon$, with the lift growing as
$\varepsilon$ tightens. The stronger statement is in the
\emph{accounting} of Table~\ref{tab:privacy}: prior work that
reports $\varepsilon\!=\!2$ in the per-iter regime is actually
operating at user-level $\varepsilon\!=\!812$, a $400\!\times\!$
overclaim that sample-once eliminates without sacrificing utility.

% ---------------------------------------------------------------------
\subsection{Q2 — Mechanism comparison: PATE wins, CANAL bounded by $R^{1/4}$}
\label{ssec:q2-mechanism}

Table~\ref{tab:mech} compares five mechanisms on the same sample-once
framework with paired student-init testing. The first two rows isolate
the CANAL closed-form allocation against the uniform-Gaussian baseline
at $K\!=\!1$ (single teacher). Rows three and four test the proposed
PATE aggregation at $K\!=\!5$ with both noise shapes. Row five reports
the best $K$ from the saturation sweep of \S\ref{ssec:q3-k-sweep}.

\begin{table}[t]
\centering
\small
\caption{Q2: mechanism comparison on DRIVE, $5$ seeds, sample-once
framework. Paired lifts in the right block are per-seed differences
against the K=1 uniform baseline (same student initialisation, same
noise precompute seed), with standard error of the mean and the
corresponding $t$-statistic. Non-private reference Dice: $0.676$. The
bold row is our headline result. \textit{CANAL} denotes the channel-WF
allocation of Theorem~\ref{thm:wf}; PATE denotes the multi-teacher
aggregation of Theorem~\ref{thm:pate-sens}.}
\label{tab:mech}
\begin{tabular}{l l l l c c c}
\toprule
   & \multicolumn{3}{c}{vessel Dice (mean $\pm$ std)}
   & \multicolumn{3}{c}{paired lift over K=1 uniform ($t$)} \\
\cmidrule(lr){2-4}\cmidrule(lr){5-7}
mechanism
   & $\varepsilon\!=\!2$ & $\varepsilon\!=\!8$ & $\varepsilon\!=\!16$
   & $\varepsilon\!=\!2$ & $\varepsilon\!=\!8$ & $\varepsilon\!=\!16$ \\
\midrule
K=1 uniform (baseline)
   & $0.581\!\pm\!0.023$ & $0.594\!\pm\!0.017$ & $0.596\!\pm\!0.013$
   & --- & --- & --- \\
K=1 CANAL
   & $0.584\!\pm\!0.025$ & $0.593\!\pm\!0.020$ & $0.599\!\pm\!0.019$
   & $+0.003\,(\!+\!1.8\sigma)$ & $-0.000\,(\!-\!0.2\sigma)$ & $+0.003\,(\!+\!1.0\sigma)$ \\
PATE K=5 + uniform
   & $0.593\!\pm\!0.016$ & $0.618\!\pm\!0.010$ & $0.627\!\pm\!0.007$
   & $+0.012\,(\!+\!2.9\sigma)$ & $+0.025\,(\!+\!5.2\sigma)$ & $+0.031\,(\!+\!7.5\sigma)$ \\
PATE K=5 + CANAL
   & $0.592\!\pm\!0.016$ & $0.619\!\pm\!0.007$ & $0.625\!\pm\!0.007$
   & $+0.011\,(\!+\!3.0\sigma)$ & $+0.025\,(\!+\!4.2\sigma)$ & $+0.030\,(\!+\!6.3\sigma)$ \\
\midrule
\textbf{PATE K=10 + uniform}
   & $\mathbf{0.607\!\pm\!0.015}$ & $\mathbf{0.628\!\pm\!0.008}$ & $\mathbf{0.648\!\pm\!0.009}$
   & $\mathbf{+0.026\,(\!+\!5.3\sigma)}$
   & $\mathbf{+0.034\,(\!+\!8.0\sigma)}$
   & $\mathbf{+0.052\,(\!+\!15.3\sigma)}$ \\
\bottomrule
\end{tabular}
\end{table}

\paragraph{Three observations.}
\emph{(i)} \textbf{CANAL alone is empirically null on RGB fundus.} The
paired CANAL$-$uniform Dice differences at $K\!=\!1$ are within the
$\pm 0.003$ seed-noise band at every $\varepsilon$ ($|t|\!<\!2\sigma$
all three columns). This is the empirical manifestation of
Corollary~\ref{cor:r14}: with the importance ratio
$R\!\approx\!2$ on bottleneck-feature gradients, the $\sigma$ ratio
between most and least important channels is $R^{1/4}\!\approx\!1.19$,
below the Dice sensitivity of a trained decoder. CANAL remains the
optimal allocation of Theorem~\ref{thm:wf} in the importance-weighted
distortion sense; what Table~\ref{tab:mech} establishes is that its
empirical room on RGB fundus is exhausted by the $R^{1/4}$ bound.

\emph{(ii)} \textbf{PATE delivers significant Dice lift, monotonic in
$K$.} At $K\!=\!5$, PATE+uniform improves over the $K\!=\!1$ baseline by
$+0.012$ to $+0.031$ Dice ($2.9\sigma$--$7.5\sigma$) across
$\varepsilon\!\in\!\{2,8,16\}$. The headline is the bold row: at
$K\!=\!10$, the lift grows to $+0.026$ -- $+0.052$
($5.3\sigma$--$15.3\sigma$), and at $\varepsilon\!=\!16$ the distilled
student's Dice ($0.648$) \emph{exceeds} the single-teacher non-private
upper bound on the local TinyUNet ($0.620$ teacher-clean), an
ensemble-regularisation effect that is consistent with the
diversification of $K\!=\!10$ disjoint cohorts.

\emph{(iii)} \textbf{CANAL composes with PATE but does not lift it on
DRIVE.} The PATE K=5+CANAL row gives Dice within $\pm 0.0015$ of the
PATE K=5+uniform row at every $\varepsilon$, with $|t|\!<\!1.5\sigma$ on
the paired test. This is exactly the prediction of
Corollary~\ref{cor:r14}: the $R^{1/4}$ bound is independent of $K$, so
sensitivity reduction by PATE cannot create room for CANAL allocation
on a dataset where the importance spectrum is already flat. We expect
this composition to surface on data where $R$ is one to two orders of
magnitude larger; multi-modal MRI is the natural next test.

% ---------------------------------------------------------------------
\subsection{Q3 — $K$-saturation: lift grows monotonically through $K\!=\!10$}
\label{ssec:q3-k-sweep}

Theorem~\ref{thm:pate-sens} predicts $\sigma\!\propto\!1/K$. We sweep
$K\!\in\!\{1, 3, 5, 8, 10\}$ at three $\varepsilon$ values with
$5$ seeds per cell (Figure~\ref{fig:k-saturation},
Table~\ref{tab:k-sat}). Lift grows monotonically in $K$ at every
$\varepsilon$ and shows no observed saturation through $K\!=\!10$.

\begin{table}[t]
\centering
\small
\caption{Q3: PATE multi-teacher $K$-sweep. Paired lift over $K\!=\!1$
baseline; mean $\pm$ SEM of $5$ seeds, with $t$-statistic. All non-trivial
rows are highly significant. $\sigma_{\mathrm{uniform}}$ at
$\varepsilon\!=\!16$ shows the linear-in-$1/K$ scaling.}
\label{tab:k-sat}
\begin{tabular}{r r r l l l}
\toprule
   $K$ & $\Delta\!=\!2/K$
       & $\sigma$ at $\varepsilon\!=\!16$
       & lift at $\varepsilon\!=\!2$
       & lift at $\varepsilon\!=\!8$
       & lift at $\varepsilon\!=\!16$ \\
\midrule
$1$   & $2.000$ & $8.64$  & baseline $0.581$  & baseline $0.594$  & baseline $0.596$ \\
$3$   & $0.667$ & $2.88$  & $+0.008\!\pm\!0.004\,(1.7\sigma)$  & $+0.016\!\pm\!0.004\,(3.6\sigma)$  & $+0.020\!\pm\!0.003\,(6.0\sigma)$ \\
$5$   & $0.400$ & $1.73$  & $+0.012\!\pm\!0.004\,(2.9\sigma)$  & $+0.025\!\pm\!0.005\,(5.2\sigma)$  & $+0.031\!\pm\!0.004\,(7.5\sigma)$ \\
$8$   & $0.250$ & $1.08$  & $+0.020\!\pm\!0.004\,(5.0\sigma)$  & $+0.031\!\pm\!0.007\,(4.7\sigma)$  & $+0.042\!\pm\!0.005\,(8.1\sigma)$ \\
$10$  & $0.200$ & $0.86$  & $\mathbf{+0.026\!\pm\!0.005\,(5.3\sigma)}$
                        & $\mathbf{+0.034\!\pm\!0.004\,(8.0\sigma)}$
                        & $\mathbf{+0.052\!\pm\!0.003\,(15.3\sigma)}$ \\
\bottomrule
\end{tabular}
\end{table}

\begin{figure}[t]
\centering
[\textit{K-saturation figure placeholder}]
\caption{Q3: paired vessel-Dice lift over the $K\!=\!1$ baseline as a
function of the PATE ensemble size $K$ (5 seeds per cell). The Dice
lift is monotonic in $K$ at all three $\varepsilon$ values and shows
no saturation through $K\!=\!10$. The vertical separation between the
three curves reflects the better resilience of larger $\varepsilon$ to
PATE-induced noise reduction (smaller $\sigma$ matters less when
$\sigma$ is already small). Source:
\texttt{drive\_pate\_K\_saturation.png}.}
\label{fig:k-saturation}
\end{figure}

\paragraph{Cohort-size caveat.}
At $K\!=\!10$ each teacher trains on $2$ DRIVE patients. The
monotonicity through $K\!=\!10$ suggests the ensemble effect dominates
per-teacher under-training in this regime, but extrapolation to
$K\!>\!10$ would require more patients than DRIVE provides. In a
deployed federated setting each ``teacher'' is one hospital with
hundreds of patients; we report the DRIVE numbers as a controlled-
setting validation, with the understanding that the lift in
Table~\ref{tab:k-sat} is a \emph{lower bound} on what a real multi-
hospital deployment would achieve.

% ---------------------------------------------------------------------
\subsection{Limitations and scope}
\label{ssec:limitations}

\paragraph{Single dataset.} All numbers are reported on DRIVE; the
$R^{1/4}$ analysis of Corollary~\ref{cor:r14} predicts that CANAL's
empirical advantage surfaces on data with a sharper bottleneck-feature
importance spectrum (multi-modal MRI), and a feature-PATE story
generalises straightforwardly to any teacher--student feature-
distillation setting. We leave multi-dataset validation to future work.

\paragraph{Empirical privacy evaluation.} We report formal user-level
$\varepsilon$ via the zCDP accountant of Section~\ref{ssec:accounting}.
A complementary empirical evaluation by membership-inference and
feature-inversion attacks (\textit{e.g.,} LiRA \cite{carlini2022lira},
inversion \cite{dosovitskiy2016inverting}) would tighten the privacy
claim from a worst-case bound to an attack-resistant operating point.
We treat this evaluation as the principal future direction for a
follow-up.

\paragraph{Spatial allocation.} Theorem~\ref{thm:wf}'s closed form is
agnostic to the allocation index set: replacing the channel index by a
joint (channel, spatial) index leaves the formula unchanged. A
spatial-axis variant requires a data-independent saliency map, which
we sketch as a parallel future direction (\textit{e.g.,} a public-proxy
Frangi vesselness map on HRF retinal images). The local proof of
concept under a synthetic radial-bias saliency was negative, but is not
informative about the real-saliency setting; we report it only for
methodological transparency in the released code.

% ---------------------------------------------------------------------
%  BibTeX entries to add to references.bib (paste into your .bib):
%
%  @inproceedings{carlini2022lira,
%    title={Membership Inference Attacks From First Principles},
%    author={Carlini, Nicholas and others},
%    booktitle={IEEE Symposium on Security and Privacy (S\&P)},
%    year={2022},
%  }
%
%  @inproceedings{dosovitskiy2016inverting,
%    title={Inverting Visual Representations with Convolutional Networks},
%    author={Dosovitskiy, Alexey and Brox, Thomas},
%    booktitle={IEEE Conference on Computer Vision and Pattern Recognition (CVPR)},
%    year={2016},
%  }
%
%  @inproceedings{staal2004drive,
%    title={Ridge-Based Vessel Segmentation in Color Images of the Retina},
%    author={Staal, Joes and Abr\`amoff, Michael D. and Niemeijer, Meindert
%            and Viergever, Max A. and van Ginneken, Bram},
%    journal={IEEE Transactions on Medical Imaging},
%    year={2004},
%  }
%
%  @inproceedings{ronneberger2015unet,
%    title={U-Net: Convolutional Networks for Biomedical Image
%           Segmentation},
%    author={Ronneberger, Olaf and Fischer, Philipp and Brox, Thomas},
%    booktitle={Medical Image Computing and Computer-Assisted
%               Intervention (MICCAI)},
%    year={2015},
%  }
% ---------------------------------------------------------------------
